\documentclass[11pt,a4paper,openany]{book}

\usepackage[margin=1in]{geometry}
\usepackage[T1]{fontenc}
\usepackage[utf8]{inputenc}
\usepackage{lmodern}
\usepackage{microtype}
\usepackage{setspace}
\onehalfspacing
\usepackage{amsmath,amssymb}
\usepackage{booktabs,longtable,array,tabularx}
\usepackage{enumitem}
\usepackage{hyperref}
\usepackage{xcolor}
\usepackage{listings}

\hypersetup{
  colorlinks=true,
  linkcolor=blue!60!black,
  citecolor=green!40!black,
  urlcolor=blue!70!black
}

\lstset{
  basicstyle=\ttfamily\small,
  breaklines=true,
  frame=single,
  columns=fullflexible,
  showstringspaces=false
}

\pagestyle{plain}

\title{\Huge\textbf{Research Assistant}\\[8pt]
\Large Individual Local Release Report, User Manual, and Future Extension Plan}
\author{Prepared for Release Review}
\date{\today}

\begin{document}
\frontmatter
\maketitle
\tableofcontents
\mainmatter

\chapter{Executive Summary}

\section{Current release target}

The current `research-assistant' release is an individual local research tool.
It is designed for one researcher working in a private filesystem workspace,
with offline-by-default workflows and Git-based sharing between researchers.

The release is not a shared department server, not a database-backed
collaboration system, not an SSO/RBAC platform, and not a hosted user
interface. Those capabilities are future extensions.

\section{What has been built}

The implemented tool provides:
\begin{itemize}[nosep]
  \item local workspace initialization, validation, repair, and configuration;
  \item privacy and offline-readiness reporting;
  \item demo setup and demo run workflows;
  \item paper ingestion from arXiv IDs, local PDFs, DOI-like queries, and title
        or topic queries;
  \item source-first arXiv LaTeX inspection when source is available;
  \item parser-tool readiness diagnostics and fixture-only parser benchmark
        smoke checks;
  \item review notes, derivation worksheets, experiment records, synthesis
        proposals, traceability reports, and governance scaffolds;
  \item backup, inspect, and restore workflows with dry-run and confirmation
        safeguards;
  \item repository hygiene checks for shareable Git repositories;
  \item workspace merge dry-run, explicit apply, provenance preservation, and
        post-merge rebuild;
  \item individual Git release validation records, validation reports, fixture
        rehearsals, performance rehearsals, and release gates;
  \item release reports and release artifact manifests.
\end{itemize}

\section{Release status}

Local validation supports a limited individual pilot. Broad individual or
Git-shared release remains blocked until the following manual gates are
completed:
\begin{itemize}[nosep]
  \item real fresh-reader onboarding from the documentation;
  \item real macOS validation;
  \item real minimal parser-tool machine validation;
  \item explicit release-owner approval for tag creation;
  \item explicit release-owner approval for artifact publication.
\end{itemize}

The release gate intentionally reports those blockers. They must not be waived
or marked as passed unless the real evidence exists.

\chapter{Product Model}

\section{One researcher at a time}

The core product assumption is simple: one researcher runs the tool locally on a
private workspace. The workspace is a directory containing configuration,
papers, source records, summaries, reviews, derived artifacts, indexes, backup
archives, and governance records.

This makes the tool easy to inspect, easy to back up, and easy to debug. It
also avoids requiring a database administrator, service operator, identity
provider, or shared hosting environment for the current release.

\section{Git as the sharing layer}

Sharing is done by Git. A researcher may publish or exchange a repository that
contains only shareable artifacts. Another researcher can clone that repository,
inspect it, run repository hygiene checks, dry-run a workspace merge, apply a
safe subset only with explicit confirmation, and rebuild derived artifacts.

Git is not treated as a transactional database. Merge conflicts and research
disagreements remain human review events.

\section{Trust boundary}

Generated outputs are review material. Parser outputs, benchmark results,
derivation worksheets, experiment records, synthesis proposals, traceability
reports, readiness reports, and validation records do not certify mathematical
correctness, scientific parser accuracy, or research approval.

Accepted human conclusions must stay explicit. The tool must not silently
promote another researcher's generated artifact into accepted local conclusions.

\chapter{Workspace And Artifact Model}

\section{Local directories}

A typical initialized workspace contains:

\begin{center}
\begin{tabularx}{\textwidth}{@{}lX@{}}
\toprule
Path & Purpose \\
\midrule
`.research-assistant/' & local configuration \\
`local\_research/papers/raw/' & private raw PDFs and paper files \\
`local\_research/papers/extracted/' & extracted text and parser output \\
`local\_research/papers/source/' & cached source bundles and source records \\
`local\_research/metadata/' & metadata records \\
`local\_research/summaries/' & paper summary records \\
`local\_research/reviews/' & review records \\
`local\_research/analysis/' & derivations, synthesis, traceability, reports \\
`local\_research/experiments/' & experiment plans and run records \\
`local\_research/governance/' & validation, release, model policy, and gate records \\
`local\_research/indices/' & rebuildable indexes \\
`local\_research/exports/' & exports and backups \\
\bottomrule
\end{tabularx}
\end{center}

\section{Shareable and private files}

The release intentionally separates shareable research artifacts from private
or generated local state.

Shareable examples include sanitized summaries, metadata, source records,
links, reviews, derivation worksheets, experiment records, traceability records,
synthesis proposals, benchmark manifests, model policies, and citation graph
reports.

Private or non-shareable examples include raw PDFs, extracted paper text,
backup archives, credentials, provider keys, tokens, `.codex', `.claude',
caches, bytecode, `build/', and `dist/'.

\chapter{Command Manual}

\section{Install and initialize}

\begin{lstlisting}[language=bash]
ra --help
ra version
ra init
ra doctor
ra doctor --matrix
\end{lstlisting}

`ra init' creates an idempotent local workspace. `ra doctor' reports optional
tool status, offline status, platform status, and workflow readiness.

\section{Demo workflow}

\begin{lstlisting}[language=bash]
ra --root /tmp/research-assistant-demo init
ra --root /tmp/research-assistant-demo demo setup
ra --root /tmp/research-assistant-demo demo run
ra --root /tmp/research-assistant-demo release-report
\end{lstlisting}

The demo creates a small synthetic research workflow with review material,
derivation, experiment, traceability, governance, and backup artifacts.

\section{Source and parser inspection}

\begin{lstlisting}[language=bash]
ra source-fetch --arxiv-id 2401.00001
ra source-show --paper-id paper_example
ra source-sections --paper-id paper_example
ra source-equations --paper-id paper_example
ra source-theorems --paper-id paper_example
ra parser-tool-matrix
ra parser-benchmark-smoke
ra parse-pdf --pdf /path/to/paper.pdf
\end{lstlisting}

Source inspection is strongest when arXiv LaTeX source is available. PDF parser
workflows are useful for inspection but parser scientific accuracy is not
certified.

\section{Review and research artifacts}

\begin{lstlisting}[language=bash]
ra find --query "transport maps"
ra show --paper-id paper_example
ra review-list
ra review-show --paper-id paper_example
ra review-mark --paper-id paper_example --status approved
ra derivation create --paper-id paper_example --title "Derivation worksheet"
ra experiment create --paper-id paper_example --claim-id claim_1 --checklist-id reproducibility
ra synthesis propose --paper-id paper_example
ra traceability build --paper-id paper_example
\end{lstlisting}

These commands create or inspect review material. They do not by themselves
create accepted mathematical conclusions.

\section{Backup and restore}

\begin{lstlisting}[language=bash]
ra backup create
ra backup inspect --path /path/to/backup.tar.gz
ra --root /tmp/restore-check backup restore --path /path/to/backup.tar.gz
ra --root /tmp/restore-check backup restore --path /path/to/backup.tar.gz --no-dry-run --confirm-restore
\end{lstlisting}

Restore is dry-run by default. A real restore requires confirmation. Overwrite
requires an additional opt-in.

\chapter{Git Sharing Workflow}

\section{Repository hygiene}

\begin{lstlisting}[language=bash]
ra repository-hygiene policy
ra repository-hygiene classify local_research/summaries/example.json
ra repository-hygiene check --strict
\end{lstlisting}

Strict hygiene checks Git state, private/generated path classes, and
secret-like payload fields. It is intended to run before publishing or sharing a
research repository.

\section{Merge dry-run and apply}

\begin{lstlisting}[language=bash]
ra --root /path/to/my/workspace workspace merge --source /path/to/other/repo --dry-run
ra --root /path/to/my/workspace workspace merge --source /path/to/other/repo --apply --confirm-merge
ra --root /path/to/my/workspace workspace rebuild-derived
\end{lstlisting}

Dry-run is the default. Apply refuses to proceed without explicit confirmation,
and conflicts involving accepted review material remain blockers.

\section{Provenance}

Applied imports preserve provenance showing the source repository, source Git
commit when available, merge report ID, and merge timestamp. This makes later
review possible without pretending that imported material is locally accepted.

\chapter{Nontrivial Showcases}

\section{Limited-pilot release packet}

\begin{lstlisting}[language=bash]
ra --root /tmp/research-assistant-final-release init
ra --root /tmp/research-assistant-final-release release-report
ra --root /tmp/research-assistant-final-release repository-hygiene check --strict
ra --root /tmp/research-assistant-final-release individual-git-release validation-substitutes
ra --root /tmp/research-assistant-final-release individual-git-release fixture-rehearsal
ra --root /tmp/research-assistant-final-release individual-git-release performance --tier synthetic_git_1000 --synthetic-count 1000 --timeout-seconds 600
ra --root /tmp/research-assistant-final-release individual-git-release validation-report
ra --root /tmp/research-assistant-final-release individual-git-release gate-build
\end{lstlisting}

This demonstrates that local fixture evidence can be complete while broad
release remains blocked for real external validation and approval.

\section{Clean install from exact wheel}

\begin{lstlisting}[language=bash]
scripts/build_release_artifacts.sh
WHEEL_PATH=dist/research_assistant-0.1.0-py3-none-any.whl scripts/run_clean_install_smoke.sh
\end{lstlisting}

The explicit `WHEEL\_PATH' avoids accidentally installing whichever wheel sorts
last in `dist/'.

\section{Parser degradation inspection}

\begin{lstlisting}[language=bash]
ra doctor --matrix
ra parser-tool-matrix
ra parser-benchmark-smoke
\end{lstlisting}

These commands show what parser workflows are available on the local machine.
They are availability and fixture checks, not scientific parser certification.

\chapter{Validation Evidence}

\section{Local tests}

The release validation packet includes:
\begin{itemize}[nosep]
  \item focused individual release integration tests;
  \item industrial scaffold regression tests;
  \item `scripts/run\_fast\_tests.sh';
  \item `scripts/run\_bounded\_tests.sh';
  \item packaging and artifact build smoke;
  \item explicit wheel clean-install smoke;
  \item individual Git release gate script;
  \item clean-checkout gate reproduction.
\end{itemize}

\section{Representative performance}

The current local evidence includes synthetic Git-sharing performance through
the `synthetic\_git\_1000' tier. This validates mechanics at that synthetic
size. It does not certify performance on private real corpora.

\section{Known blocked gates}

The release gate remains blocked for:
\begin{itemize}[nosep]
  \item fresh-reader onboarding not yet recorded;
  \item macOS validation not yet recorded;
  \item minimal parser-tool machine validation not yet recorded;
  \item release-owner tag approval not yet recorded;
  \item release-owner publication approval not yet recorded.
\end{itemize}

\chapter{Maintainer Notes}

\section{Where to look}

Maintainers should start with:
\begin{itemize}[nosep]
  \item `docs/maintainer\_guide.md';
  \item `docs/proposal/individual\_git\_release\_target.md';
  \item `docs/release\_checklist.md';
  \item `src/research\_assistant/cli.py';
  \item `src/research\_assistant/individual\_release.py';
  \item `src/research\_assistant/individual\_git\_release.py'.
\end{itemize}

\section{High-risk changes}

Treat these changes as release-critical:
\begin{itemize}[nosep]
  \item changing backup restore confirmation semantics;
  \item changing merge dry-run or apply confirmation behavior;
  \item changing forbidden path or secret scanning policy;
  \item changing validation scopes or gate readiness flags;
  \item weakening parser limitation language;
  \item changing JSON schema versions or artifact IDs.
\end{itemize}

\chapter{Known Limitations}

\begin{itemize}[nosep]
  \item The current release is local and file-based.
  \item Git sharing is not real-time collaboration.
  \item Parser scientific accuracy is not certified.
  \item Generated artifacts remain review material.
  \item Real external platform and human onboarding evidence is still missing.
  \item Tagging and publication require explicit release-owner approval.
  \item Multi-user database/service/RBAC/hosted UI work is future scope.
\end{itemize}

\chapter{Future Multi-User Extension}

The future extension path may include:
\begin{itemize}[nosep]
  \item shared storage with migration policy;
  \item service deployment and background jobs;
  \item SSO/RBAC and audit-log policy;
  \item hosted UI for inspection and review;
  \item department operations and support SOPs;
  \item stronger parser/source benchmarks on gold corpora;
  \item real-world scalability validation on sanitized corpora;
  \item governed live-provider or LLM integration.
\end{itemize}

Those are not prerequisites for the current individual release. They should be
planned and validated separately so the local tool can remain simple,
inspectable, and useful.

\backmatter
\chapter{Release Checklist Summary}

\begin{lstlisting}[language=bash]
scripts/run_fast_tests.sh
scripts/run_bounded_tests.sh
PYTHONPATH=src python -m pytest tests/integration/test_individual_release_cli.py -q
scripts/build_release_artifacts.sh
WHEEL_PATH=dist/research_assistant-0.1.0-py3-none-any.whl scripts/run_clean_install_smoke.sh
scripts/run_individual_git_release_gate.sh
git diff --check
git status --short --ignored
\end{lstlisting}

Final broad release requires real external validation and release-owner
approval. Until then, the tool is suitable for limited individual pilot use with
documented limitations.

\end{document}
